\section{Deleting root evaluation from the raw clause calculus}
Let $\mathsf{CE5}_3$ denote the pinned raw clause calculus, whose equality
operation has the three rows Root, Zero, and Successor.  Let
$\mathsf{CE5}_2$ be exactly the same calculus with Root omitted.  In
particular, terms, formulas, clause contexts, the designated contradiction,
all fresh-variable conditions, and ordinary capture-avoiding substitution
are unchanged.  The zero and successor rows are
\[
 \frac{\Gamma\vdash s=t\qquad
       \Gamma\vdash I(0,s;u.h)=r}
      {\Gamma\vdash r=t}\;(E_0),
\qquad
 \frac{\Gamma\vdash I(n,a;u.h)=t\qquad
       \Gamma\vdash I(Sn,a;u.h)=r}
      {\Gamma\vdash r=h[t/u]}\;(E_S).
\]
All other shapes are rejected.  In the standalone implementation these
are two branches of one deterministic partial operation, not proof search.

\paragraph{Identity-frame lemma.}
Suppose reflexivity is available, and a term $W(s)$ supports
$s=t,W(s)=r\vdash r=t$.  Then $s=t,s=r\vdash r=t$ is derivable without
assuming a computation equation $W(s)=s$.  From $q:s=r$ and
$\operatorname{Refl}(W(s))$, the displayed rule gives $W(s)=r$;
from $p:s=t$ and that result it gives $r=t$.

For the raw iterator use $W(s)=I(0,s;u.0)$.  It has exactly the free
variables of $s$, binds no variable occurring in the inserted $s$, and
is well scoped wherever $s$ is well scoped.  The derived proof is
\[
 \operatorname{RC}(p,q)=
 E_0\bigl(p,E_0(q,\operatorname{Refl}(W(s)))\bigr).
\]
The first zero evaluation has source $s$, target $r$, minor left side
$W(s)$, and minor right side $W(s)$, hence conclusion $W(s)=r$.
The second has source $s$, target $t$, and minor right side $r$,
hence conclusion $r=t$.  The update is the closed zero term; its bound
variable never occurs.  Neither induction, logical introduction,
normalization, congruence, nor a numerical-substitution theorem is
used by this construction.

\paragraph{Exact consequence theorem.}
For every finite well-scoped clause context $\Gamma$ and Eq/$D$ formula
$F$, one has
\[
 \Gamma\vdash_{\mathsf{CE5}_3}F
 \quad\Longleftrightarrow\quad
 \Gamma\vdash_{\mathsf{CE5}_2}F.
\]
The reverse implication is inclusion.  For the forward implication,
induct on the finite proof tree.  Copy every reflexivity, clause-use,
paired-elimination and paired-introduction node, recursively translating
its proof premises.  Copy the Zero and Successor evaluation nodes.
Replace each Root node by the identity-frame derivation above.

Every recursive replacement has exactly the same conclusion in exactly
the same clause context and rigid-variable scope as the old premise.
The fresh eigenvariables at logical nodes are retained, not substituted.
The matrices and the witnesses of paired introduction remain literally
unchanged, so the equation-root/literal-zero mode test is unchanged.
There is therefore no hidden assumption of direct substitution
admissibility for the witness-selected introduction rule.  Arbitrary
initial reusable clauses are allowed; the construction does not depend
on how they entered the context.

\paragraph{Proof size.}
If the old proof has $N$ inference nodes and contains $k$ Root nodes,
the transformed proof has exactly $N+2k$ nodes.  Each replaced node
becomes two Zero nodes and one reflexivity leaf, with each of the
original premise subtrees occurring once.  All other proof nodes are
copied.  The bound counts tree nodes, not bytes: $W(s)$ copies the
source term, so there is no claimed constant certificate-byte overhead.

\paragraph{Why the elementary frame proof is already short.}
It adds exactly one reflexivity leaf and two Zero evaluations per
Root.  For generic equations $x=y$ and $x=z$, with $x,y,z$ distinct
variables, neither supplied equation has an iterator as left side.
Consequently neither can serve as the minor premise of the first
Zero or Successor evaluation.  A derivation using only these two
inputs and evaluation, with no further equality leaf, cannot even
start.  With one added reflexivity leaf, a single Zero evaluation
cannot yield $z=y$: its iterator-headed minor must either be the
reflexivity leaf, yielding an iterator on the conclusion's left side,
or be a previously derived equation, which requires an earlier step.
Thus the one-leaf/two-step simulation is locally minimal in this
restricted direct equality fragment.  No global shortest-proof or
axiomatization lower bound is asserted.

\paragraph{Bi-interpretability.}
The new and old clause theories use identical numerical terms and
identical Eq/$D$ formulas.  The proof translations leave all of them
unchanged.  Their mutual interpretations are therefore identities,
with numerical graph $x=y$ and literal identity formula composites.
Composing with the pinned raw-clause-to-native bridge and the native
whole-domain, actual-equality interpretation with iPRA transfers both
formula round trips and the numerical identity graphs.  This is
stronger than a mere equivalence of consistency claims.  In particular,
consistency for the unchanged contradiction $0=S0$ is equivalent.

\paragraph{Dependence and scope.}
The identity-frame lemma itself is inherited from
\texttt{equality\_from\_iteration}; the clause representation is inherited
from \texttt{raw\_clausal}.  The new result here is their exact combination,
with a controlled complete-checker comparison and an explicit proof
translator.  It is a modest simplification, not a claim of a new arithmetic
principle, historical priority, absolute minimality, or a new global
compact-code record.
